% tw_cruciform.tex -- Chapter 5: the cruciform section and its torsional mode.
\begin{tikzpicture}[>=Stealth,scale=1.0]
  % left: the section
  \begin{scope}
    \fill[cu!12] (-1.5,-0.09) rectangle (1.5,0.09);
    \fill[cu!12] (-0.09,-1.5) rectangle (0.09,1.5);
    \draw[cu,line width=1.2pt] (-1.5,0) -- (1.5,0) (0,-1.5) -- (0,1.5);
    \fill (0,0) circle (1.6pt);
    \node[font=\small,below left] at (0,0) {$C = S$};
    \draw[<->] (0.12,0.35) -- (1.5,0.35) node[midway,above] {\small $b$};
    \node[font=\small,right] at (1.5,0) {$t$};
    \draw[->,orange!80!black,thick] (1.1,1.0) arc (40:80:1.4);
    \node[font=\small,orange!80!black] at (1.55,1.35) {$\varphi$};
    \node[font=\small,align=center] at (0,-2.0) {four legs meet at one point:\\ no warping, $\Gamma \approx 0$};
  \end{scope}
  % right: the section seen from above at five heights, rotated by phi(x) = phi_max sin(pi x/L)
  \begin{scope}[shift={(3.4,0)}]
    \foreach \k/\lab in {0/0,1/{L/4},2/{L/2},3/{3L/4},4/L} {
      \pgfmathsetmacro{\a}{35*sin(45*\k)}
      \begin{scope}[shift={(1.25*\k,0)},rotate=\a]
        \draw[gray!50] (-0.5,0) -- (0.5,0) (0,-0.5) -- (0,0.5);
      \end{scope}
      \begin{scope}[shift={(1.25*\k,0)},rotate=\a]
        \draw[cu,line width=1.1pt] (-0.5,0) -- (0.5,0) (0,-0.5) -- (0,0.5);
      \end{scope}
      \node[font=\small] at (1.25*\k,-0.9) {$\lab$};
    }
    \node[font=\small,left] at (-0.6,-0.9) {$x$:};
    \draw[->,orange!80!black,thick] (2.5,0.8) arc (90:125:0.8);
    \node[font=\small,orange!80!black] at (2.75,1.0) {$\varphi$};
    \node[font=\small,align=center] at (2.5,-1.75) {the axis stays straight, the section rotates:\\
      $\sigma = G(t/b)^2$ for any length};
  \end{scope}
\end{tikzpicture}
